% فصل دوم: مبانی نظری و پیشینه پژوهش
\chapter{مبانی نظری و پیشینه پژوهش}
\label{ch:background}

هدف این فصل، تبیین و تشریح مفاهیم پایه‌ای است که چارچوب شبکه‌های عصبی آگاه از فیزیک بر مبنای آن‌ها شکل گرفته است. در ابتدا، تعاریف و ویژگی‌های کلی معادلات دیفرانسیل جزئی به‌همراه مثال‌های کلاسیک ارائه می‌گردد و سپس روش‌های عددی کلاسیک حل این معادلات به اختصار بررسی می‌شوند. در ادامه، ساختار شبکه‌های عصبی مصنوعی، الگوریتم‌های آموزش و بهینه‌سازی و تکنیک مشتق‌گیری خودکار تشریح می‌گردند. در بخش پایانی فصل، پیشینه پژوهش‌های پیشین در حوزه ترکیب یادگیری ماشین با اصول فیزیکی مرور می‌شود تا جایگاه پژوهش حاضر مشخص گردد.

\section{معادلات دیفرانسیل جزئی}
\label{sec:pde}

معادله دیفرانسیل جزئی، رابطه‌ای ریاضی میان یک تابع مجهول چندمتغیره و مشتقات جزئی آن نسبت به متغیرهای مستقل است. صورت عمومی یک معادله دیفرانسیل جزئی برای تابع مجهول $u(\mathbf{x})$ به فرم زیر بیان می‌شود:
\begin{equation}
F\left(\mathbf{x}, u, Du, D^2u, \dots, D^ku\right) = 0,
\label{eq:pde-general}
\end{equation}
که در آن $\mathbf{x} = (x_1, \dots, x_d, t) \in \mathbb{R}^{d+1}$ بردار متغیرهای مستقل مکانی و زمانی، $Du$ نمایانگر گرادیان یا مجموعه مشتقات جزئی مرتبه اول، $D^2u$ ماتریس هشین یا مشتقات جزئی مرتبه دوم و $D^ku$ مشتقات جزئی مرتبه $k$ام هستند. بالاترین مرتبه مشتق موجود در رابطه را مرتبه معادله می‌نامند. چنانچه عملگر $F$ نسبت به متغیر $u$ و کلیه مشتقات آن خطی باشد، معادله خطی و در غیر این صورت، غیرخطی\LTRfootnote{Nonlinear} نامیده می‌شود \cite{evans2010}.

برای تعیین جواب یکتا و فیزیکی از میان بی‌شمار جواب‌های عمومی معادله، مسئله باید همراه با شرایط مرزی\LTRfootnote{Boundary Conditions} و شرایط اولیه\LTRfootnote{Initial Conditions} فرمول‌بندی شود. شرط اولیه حالت سیستم را در لحظه آغازین $t=0$ مشخص نموده و شرایط مرزی نحوه برهم‌کنش سیستم با مرزهای دامنه مکانی $\partial\Omega$ را تعیین می‌کنند. بر اساس تعریف هادامارد، مسئله‌ای که دارای جواب یکتا باشد و جواب آن نسبت به داده‌های اولیه و مرزی پیوسته و پایدار بماند، یک مسئله خوش‌وضع\LTRfootnote{Well-posed Problem} نامیده می‌شود.

در محاسبات علمی، برخی معادلات به‌دلیل تبیین پدیده‌های بنیادین، به‌عنوان مسائل معیار شناخته می‌شوند:

\textbf{معادله گرما.} این معادله ساده‌ترین و بنیادی‌ترین نمونه از معادلات سهموی است که انتقال و پخش دما را توصیف می‌کند. فرم یک‌بعدی آن عبارت است از:
\begin{equation}
u_t = \alpha u_{xx},
\label{eq:heat}
\end{equation}
که در آن $u(t,x)$ بیانگر توزیع دما در زمان $t$ و مکان $x$، و $\alpha > 0$ ضریب نفوذ یا پخش حرارتی است.

\textbf{معادله برگرز.} معادله برگرز\LTRfootnote{Burgers Equation} نمونه‌ای بارز از معادلات دیفرانسیل جزئی غیرخطی است که می‌توان آن را مدل تقلیل‌یافته یک‌بعدی از معادلات ناویر-استوکس در دینامیک سیالات و مدل‌های جریان ترافیک دانست \cite{hopf1950}. صورت استاندارد معادله برگرز با لزجت به فرم زیر نوشته می‌شود:
\begin{equation}
u_t + u u_x - \nu u_{xx} = 0,
\label{eq:burgers-intro}
\end{equation}
که در آن $\nu > 0$ ضریب لزجت سینماتیکی است. در این معادله، جمله غیرخطی $u u_x$ نمایانگر انتقال (همرفت) است که تمایل به تیز کردن پروفیل موج و ایجاد ناپیوستگی دارد، در حالی که جمله $\nu u_{xx}$ اثر پخشی داشته و به هموارسازی گرادیان‌ها می‌انجامد. تعامل این دو جمله، به‌ویژه در مقادیر بسیار کوچک $\nu$، منجر به پدیدار شدن جبهه‌های شیب‌دار تند و شبیه به موج شوک می‌شود. به همین علت، این معادله به معیاری چالش‌برانگیز و معتبر برای سنجش توانمندی روش‌های عددی تبدیل شده است و در این پژوهش نیز به‌عنوان مسئله آزمون مورد استفاده قرار گرفته است.

\textbf{معادله غیرخطی شرودینگر.} این معادله در اپتیک غیرخطی، فیزیک پلاسما و مکانیک کوانتومی کاربرد گسترده‌ای دارد و در مقاله مرجع \cite{raissi2019} نیز به‌عنوان یکی از مسائل آزمون بررسی گردیده است.

\section{روش‌های عددی کلاسیک}
\label{sec:classical}

به‌جز در موارد خاص و با هندسه‌های بسیار متقارن، معادلات دیفرانسیل جزئی غیرخطی فاقد جواب تحلیلی بسته هستند؛ از این رو، روش‌های عددی نقش محوری را در حل عملی این مسائل بر عهده دارند. اساس کار روش‌های عددی، گسسته‌سازی دامنه پیوسته و تبدیل معادله دیفرانسیل به یک دستگاه معادلات جبری متناهی است \cite{leveque2007}. سه رویکرد اصلی در این زمینه عبارت‌اند از:

\textbf{روش تفاضل‌های محدود.} در روش تفاضل‌های محدود\LTRfootnote{Finite Difference Method}، دامنه پیوسته با یک شبکه منظم از نقاط گسسته جایگزین شده و مشتقات جزئی با استفاده از بسط تیلور پیرامون نقاط همسایه تقریب زده می‌شوند. به عنوان نمونه، تقریب تفاضل مرکزی مرتبه دوم برای مشتق دوم مکانی به صورت زیر است:
\begin{equation}
u_{xx}(x_i) \approx \frac{u_{i+1} - 2u_i + u_{i-1}}{(\Delta x)^2},
\label{eq:fd}
\end{equation}
که در آن $\Delta x$ طول گام گسسته‌سازی مکانی است. مزیت عمده این روش، سادگی مفاهیم و سهولت پیاده‌سازی است؛ اما در مواجهه با دامنه‌های هندسی پیچیده و ناهموار، کارایی آن به شدت کاهش می‌یابد.

\textbf{روش اجزای محدود.} در روش اجزای محدود\LTRfootnote{Finite Element Method}، دامنه مکانی به زیردامنه‌های کوچکی به نام المان (نظیر مثلث‌ها یا چهارضلعی‌ها) تقسیم می‌گردد و جواب معادله بر حسب توابع پایه چندجمله‌ای موضعی بسط داده می‌شود. سپس با اعمال فرمول‌بندی ضعیف (روش وردشی یا گالرکین)، دستگاه معادلات جبری تشکیل می‌گردد. انعطاف‌پذیری فوق‌العاده در مدیریت هندسه‌های نامنظم و شرایط مرزی ترکیبی، این روش را به استاندارد مهندسی بدل ساخته است؛ هرچند پیچیدگی محاسباتی و هزینه الگوریتمی آن نسبتاً بالاست.

\textbf{روش‌های طیفی.} در روش‌های طیفی\LTRfootnote{Spectral Methods}، تابع جواب به‌صورت ترکیبی خطی از توابع پایه متعامد سراسری (نظیر چندجمله‌ای‌های چبیشف، لژاندر یا توابع مثلثاتی فوریه) تقریب زده می‌شود. مزیت برجسته روش‌های طیفی، نرخ همگرایی نمایی (موسوم به دقت طیفی) برای توابع با جواب‌های بی‌نهایت‌بار هموار است \cite{trefethen2000}. در بخش تجربی این پایان‌نامه، از یک حل‌گر طیفی فوریه برای تولید داده‌های مرجع دقیق معادله برگرز استفاده شده است.

جدول \ref{tab:numerical} مقایسه‌ای اجمالی از ویژگی‌های این سه دسته روش را نشان می‌دهد. تمامی روش‌های کلاسیک در دو ویژگی اشتراک دارند: وابستگی ساختاری به تولید شبکه محاسباتی و ضرورت حل مجدد کل مسئله در صورت تغییر شرایط اولیه، شرایط مرزی یا پارامترهای سیستم. همچنین با افزایش ابعاد مسئله ($d \ge 3$)، پدیده نفرین ابعاد منجر به رشد نمایی تعداد درجات آزادی و هزینه‌های حافظه و محاسبات می‌گردد.

\begin{table}[t]
\caption{مقایسه ویژگی‌های روش‌های عددی کلاسیک در حل معادلات دیفرانسیل جزئی}
\label{tab:numerical}
\centering
\singlespacing
\begin{tabular}{|p{2.6cm}|p{3.6cm}|p{3.5cm}|p{3.7cm}|}
\hline
\textbf{روش محاسباتی} & \textbf{ایده بنیادین} & \textbf{مزیت اصلی} & \textbf{محدودیت اصلی} \\
\hline
تفاضل‌های محدود & تقریب موضعی مشتقات روی شبکه ساخت‌یافته & سادگی الگوریتمی و پیاده‌سازی & ناتوانی در هندسه‌های نامنظم \\
\hline
اجزای محدود & فرمول‌بندی ضعیف روی المان‌های نامنظم & انعطاف‌پذیری هندسی بالا & هزینه بالای تولید شبکه و محاسبات \\
\hline
روش‌های طیفی & بسط جواب بر حسب توابع پایه سراسری & نرخ همگرایی نمایی (دقت بسیار بالا) & محدود به هندسه‌های ساده و شرایط هموار \\
\hline
\end{tabular}
\end{table}

\section{شبکه‌های عصبی مصنوعی}
\label{sec:ann}

شبکه عصبی مصنوعی\LTRfootnote{Artificial Neural Network (ANN)} مدلی ریاضی و محاسباتی است که از پردازش اطلاعات در شبکه‌های بیولوژیکی الهام گرفته شده است. واحد بنیادی پردازش در این ساختارها، نورون مصنوعی است. هر نورون حاصل‌جمع وزن‌دار ورودی‌ها را با یک مقدار بایاس جمع نموده و آن را از یک تابع غیرخطی به نام تابع فعال‌ساز\LTRfootnote{Activation Function} عبور می‌دهد:
\begin{equation}
y = \sigma\left(\sum_{i=1}^{n} w_i x_i + b\right),
\label{eq:neuron}
\end{equation}
که در آن $x_i$ ورودی‌ها، $w_i$ وزن‌های سیناپسی، $b$ پارامتر بایاس و $\sigma(\cdot)$ تابع فعال‌ساز است. چیدمان لایه‌ای این واحدها، ساختار پرسپترون چندلایه\LTRfootnote{Multi-Layer Perceptron (MLP)} را پدید می‌آورد. تبدیل برداری در لایه $l$ام به‌صورت رابطه بازگشتی زیر تعریف می‌شود:
\begin{equation}
\mathbf{h}^{(l)} = \sigma\left(\mathbf{W}^{(l)} \mathbf{h}^{(l-1)} + \mathbf{b}^{(l)}\right),
\label{eq:mlp}
\end{equation}
که در آن $\mathbf{h}^{(0)} = \mathbf{x}$ بردار ورودی، $\mathbf{W}^{(l)}$ ماتریس وزن‌ها و $\mathbf{b}^{(l)}$ بردار بایاس لایه $l$ام است. ساختار گرافیکی یک پرسپترون چندلایه در شکل \ref{fig:mlp} به تصویر کشیده شده است.

\begin{figure}[t]
\centering
\begin{latin}
\begin{tikzpicture}[>=Stealth,
  neuron/.style={circle, draw, minimum size=7mm, inner sep=0pt},
  ann/.style={->, thin}]
  % input layer
  \node[neuron] (i1) {$x_1$};
  \node[neuron, below=0.6cm of i1] (i2) {$x_2$};
  % hidden layer 1
  \node[neuron, right=1.8cm of i1, yshift=0.6cm] (h11) {};
  \node[neuron, below=0.6cm of h11] (h12) {};
  \node[neuron, below=0.6cm of h12] (h13) {};
  % hidden layer 2
  \node[neuron, right=1.8cm of h11] (h21) {};
  \node[neuron, below=0.6cm of h21] (h22) {};
  \node[neuron, below=0.6cm of h22] (h23) {};
  % output
  \node[neuron, right=1.8cm of h22] (o1) {$y$};
  % connections
  \foreach \s in {i1,i2} \foreach \d in {h11,h12,h13} \draw[ann] (\s) -- (\d);
  \foreach \s in {h11,h12,h13} \foreach \d in {h21,h22,h23} \draw[ann] (\s) -- (\d);
  \foreach \s in {h21,h22,h23} \draw[ann] (\s) -- (o1);
  % labels
  \node[above=0.35cm of h11, xshift=0.9cm, font=\small] {Hidden layers};
  \node[above=0.35cm of i1, font=\small] {Input};
  \node[above=0.35cm of o1, font=\small] {Output};
\end{tikzpicture}
\end{latin}
\caption{ساختار شماتیک یک شبکه عصبی پرسپترون چندلایه با دو متغیر ورودی، دو لایه پنهان و یک خروجی اسکالر}
\label{fig:mlp}
\end{figure}

انتخاب تابع فعال‌ساز مناسب نقشی تعیین‌کننده در رفتار و قابلیت مشتق‌پذیری شبکه دارد. جدول \ref{tab:activations} توابع فعال‌ساز متداول را مقایسه می‌نماید. در شبکه‌های عصبی آگاه از فیزیک، به‌کارگیری توابع بی‌نهایت‌بار مشتق‌پذیر و هموار نظیر تانژانت هیپربولیک ($\tanh$) امری ضروری است؛ زیرا تشکیل جملات دیفرانسیلی نیازمند محاسبه مشتقات مرتبه دوم یا بالاتر خروجی نسبت به ورودی‌هاست، در حالی که توابعی نظیر \lr{ReLU} دارای مشتق دوم همه‌جا صفر هستند.

\begin{table}[t]
\caption{توابع فعال‌ساز متداول در شبکه‌های عصبی عمیق}
\label{tab:activations}
\centering
\singlespacing
\begin{tabular}{|p{3.8cm}|>{\centering\arraybackslash}p{4.6cm}|p{5.0cm}|}
\hline
\textbf{عنوان تابع} & \textbf{فرمول ریاضی} & \textbf{ویژگی رفتاری} \\
\hline
سیگموئید (\lr{Sigmoid}) & $\sigma(z) = \dfrac{1}{1+e^{-z}}$ & بازه خروجی $(0,1)$، هموار و مشتق‌پذیر \\
\hline
تانژانت هیپربولیک (\lr{Tanh}) & $\tanh(z) = \dfrac{e^{z}-e^{-z}}{e^{z}+e^{-z}}$ & بازه خروجی $(-1,1)$، هموار، متقارن حول مبدأ \\
\hline
واحد خطی یکسوشده (\lr{ReLU}) & $\mathrm{ReLU}(z) = \max(0,z)$ & محاسبه سریع، نامشتق‌پذیر در صفر، مشتق دوم صفر \\
\hline
\end{tabular}
\end{table}

پایه‌های نظری توانمندی شبکه‌های عصبی در تقریب توابع با «قضیه تقریب عمومی» استوار شده است. بر اساس قضیه اثبات‌شده توسط سیبنکو \cite{cybenko1989}:

\begin{theorem}[تقریب عمومی سیبنکو]
فرض کنید $\sigma(\cdot)$ یک تابع پیوسته سیگموئیدی (غیرخطی و کراندار) باشد. در این صورت، هر تابع پیوسته $f \in C(K)$ تعریف‌شده روی یک مجموعه فشرده $K \subset \mathbb{R}^n$ را می‌توان با یک شبکه عصبی پیش‌خور تک‌لایه پنهان با دقت دلخواه $\epsilon > 0$ تقریب زد.
\end{theorem}

هرچند این قضیه کفایت تئوریک شبکه‌های تک‌لایه با پهنای نامتناهی را اثبات می‌کند، اما در عمل شبکه‌های عمیق‌تر (با لایه‌های متوالی) با تعداد کل پارامترهای بسیار کمتر، قابلیت یادگیری توابع پیچیده فیزیکی را فراهم می‌سازند.

\section{آموزش شبکه‌های عصبی عمیق}
\label{sec:training}

فرایند آموزش شبکه عصبی به معنای یافتن مجموعه پارامترهای بهینه $\theta^* = \{\mathbf{W}, \mathbf{b}\}$ است به‌گونه‌ای که یک تابع هدف مشخص کمینه گردد. در مسائل رگرسیون پیوسته، معیار میانگین مربعات خطا\LTRfootnote{Mean Squared Error (MSE)} پرکاربردترین تابع هزینه است:
\begin{equation}
\mathrm{MSE}(\theta) = \frac{1}{N}\sum_{i=1}^{N} \left(u_\theta(\mathbf{x}_i) - u_i\right)^2,
\label{eq:mse}
\end{equation}
که در آن $u_\theta(\mathbf{x}_i)$ پیش‌بینی مدل برای ورودی $\mathbf{x}_i$ و $u_i$ مقدار هدف است.

کمینه‌سازی تابع هزینه با الگوریتم‌های گرادیان کاهشی انجام می‌شود:
\begin{equation}
\theta_{k+1} = \theta_k - \eta \nabla_\theta \mathcal{L}(\theta_k),
\label{eq:gd}
\end{equation}
که در آن $\eta > 0$ نرخ یادگیری\LTRfootnote{Learning Rate} و $\nabla_\theta \mathcal{L}$ گرادیان تابع هزینه نسبت به بردار پارامترهاست. محاسبه کارآمد این گرادیان‌ها توسط الگوریتم پس‌انتشار خطا\LTRfootnote{Backpropagation} صورت می‌گیرد که مبتنی بر قاعده زنجیره‌ای روی گراف محاسباتی است \cite{rumelhart1986}.

در بهینه‌سازی شبکه‌های آگاه از فیزیک، ترکیب دو بهینه‌ساز زیر راهبردی استاندارد به شمار می‌رود:

\textbf{بهینه‌ساز آدام (\lr{Adam}).} این بهینه‌ساز با محاسبه گشتاور اول و دوم متحرکِ گرادیان‌ها، نرخ یادگیری تطبیقی را برای هر پارامتر اعمال می‌کند \cite{kingma2015}:
\begin{equation}
\begin{aligned}
m_k &= \beta_1 m_{k-1} + (1-\beta_1) g_k, \\
v_k &= \beta_2 v_{k-1} + (1-\beta_2) g_k^2, \\
\hat{m}_k &= \frac{m_k}{1-\beta_1^k}, \quad \hat{v}_k = \frac{v_k}{1-\beta_2^k}, \\
\theta_{k+1} &= \theta_k - \eta \frac{\hat{m}_k}{\sqrt{\hat{v}_k} + \epsilon},
\end{aligned}
\label{eq:adam}
\end{equation}
که در آن $g_k = \nabla_\theta \mathcal{L}(\theta_k)$، $\beta_1=0.9$، $\beta_2=0.999$ و $\epsilon=10^{-8}$ مقادیر استاندارد هستند. عبارات $\hat{m}_k$ و $\hat{v}_k$ اصلاح‌کننده بایاس اولیه در گام‌های آغازین بهینه‌سازی می‌باشند.

\textbf{بهینه‌ساز شبه‌نیوتنی ال-بی‌اف‌جی‌اس (\lr{L-BFGS}).} این روش شبه‌نیوتنی با تقریب ماتریس معکوس هشین بر پایه تغییرات گرادیان در گام‌های پیشین، همگرایی مرتبه دو و فوق‌خطی در نزدیکی نقاط کمینه ایجاد می‌کند \cite{liu1989}. در آموزش شبکه‌های آگاه از فیزیک، معمولاً ابتدا چند هزار گام با آدام اجرا می‌شود تا بهینه‌سازی از نواحی ناملایم عبور کند، و سپس برای رسیدن به دقت‌های بالای مهندسی، آموزش با الگوریتم \lr{L-BFGS} تکمیل می‌گردد.

علاوه بر بهینه‌ساز، مقداردهی اولیه وزن‌ها نقشی حیاتی در همگرایی دارد. در این پژوهش از روش مقداردهی گزاویه\LTRfootnote{Xavier Initialization} استفاده شده است که واریانس وزن‌های هر لایه را متناسب با تعداد نورون‌های ورودی و خروجی آن تنظیم می‌نماید \cite{glorot2010}.

\section{مشتق‌گیری خودکار}
\label{sec:autodiff}

مشتق‌گیری خودکار\LTRfootnote{Automatic Differentiation} مجموعه‌ای از تکنیک‌های محاسباتی برای ارزیابی مشتق توابع پیاده‌سازی‌شده در کدهای برنامه‌نویسی است \cite{baydin2018}. مشتق‌گیری خودکار از دو رویکرد کلاسیک دیگر متمایز است: بر خلاف مشتق‌گیری عددی (نظیر تفاضل محدود)، فاقد خطای برش و ناپایداری ناشی از گردکردن است و تا دقت ممیز شناور ماشین دقیق عمل می‌کند؛ و بر خلاف مشتق‌گیری نمادین، با عبارات جبری پیچیده روبه‌رو نمی‌شود و از پدیده رشد نمایی طول عبارات جبری مبرا است.

دو حالت اصلی در مشتق‌گیری خودکار عبارت‌اند از:
\begin{enumerate}
\item \textbf{حالت مستقیم (\lr{Forward Mode}):} مشتقات هم‌زمان با اجرای مستقیم برنامه و با محاسبه مشتقات جهت‌دار ارزیابی می‌شوند. این حالت برای توابعی با ابعاد ورودی اندک و خروجی زیاد کارآمد است.
\item \textbf{حالت معکوس (\lr{Reverse Mode}):} ابتدا اجرای مستقیم صورت گرفته و مقادیر متغیرهای میانی ذخیره می‌شوند؛ سپس با یک پیمایش بازگشتی از خروجی به سمت ورودی‌ها، گرادیان کامل نسبت به تمامی ورودی‌ها محاسبه می‌گردد. این حالت زیربنای الگوریتم پس‌انتشار است.
\end{enumerate}

در چارچوب شبکه‌های آگاه از فیزیک، مشتق‌گیری خودکار کاربردی دوگانه دارد: نخست، محاسبه گرادیان تابع هزینه نسبت به وزن‌های شبکه جهت به‌روزرسانی پارامترها؛ و دوم (که مشخصه اختصاصی این چارچوب است)، محاسبه مشتقات جزئی متغیر وابسته نسبت به متغیرهای مستقل مکانی و زمانی ($u_t, u_x, u_{xx}$) برای ارزیابی مستقیم باقیمانده معادلات دیفرانسیل. کتابخانه‌های مدرن یادگیری عمیق نظیر جکس \cite{jax2018} و تنسورفلو \cite{abadi2016} این قابلیت را با کامپایل کارآمد روی سخت‌افزارهای پردازشی فراهم می‌سازند.

\section{پیشینه پژوهش}
\label{sec:related}

به‌کارگیری شبکه‌های عصبی در حل معادلات دیفرانسیل ریشه در پژوهش‌های دهه ۱۹۹۰ میلادی دارد. لاگاریس و همکارانش در سال ۱۹۹۸ نشان دادند که می‌توان جواب معادلات دیفرانسیل معمولی و جزئی را با پرسپترون‌های چندلایه فرمول‌بندی کرد به‌طوری که شرایط مرزی به‌صورت ساختاری ارضا شده و باقیمانده معادله روی نقاط آزمون کمینه گردد \cite{lagaris1998}. با این حال، محدودیت توان پردازشی و عدم وجود کتابخانه‌های مشتق‌گیری خودکار در آن برهه، توسعه این روش‌ها را متوقف ساخت.

با پیشرفت علوم داده، رویکردهای یادگیری فیزیک‌پایه مجدداً کانون توجه قرار گرفتند. در این راستا، مدل‌سازی مبتنی بر فرایندهای گاوسی\LTRfootnote{Gaussian Process} برای استنتاج عملگرهای خطی \cite{raissi2017a} و معادلات غیرخطی از طریق خطی‌سازی موضعی \cite{raissi2018} پیشنهاد گردید. این روش‌ها توانایی ارائه برآورد عدم‌قطعیت را دارا بودند، اما به‌دلیل فرضیات گاوسی، قابلیت تعمیم آن‌ها به مسائل به‌شدت غیرخطی با محدودیت روبه‌رو بود.

در رویکردی موازی، کارپاتنه و همکارانش چارچوب «علم داده هدایت‌شده با نظریه»\LTRfootnote{Theory-Guided Data Science} را برای ادغام اصول علمی در مدل‌های یادگیری ماشین معرفی نمودند \cite{karpatne2017}. همچنین روش‌های شناسایی سیستم‌های دینامیکی بر پایه تنک‌سازی نظیر \lr{SINDy} توسط برانتون و همکاران توسعه یافت \cite{sindy2016}؛ این روش‌ها اگرچه در کشف معادلات حاکم موفق بودند، اما به پایگاه داده‌ای جامع از توابع کاندید و مشتق‌گیری عددی روی داده‌های نویزی وابسته بودند.

انتشار مقاله رئیسی، پردیکاریس و کارنیاداکیس در سال ۲۰۱۹ تحولی یکپارچه در این حوزه به وجود آورد \cite{raissi2019}. آن‌ها با ترکیب شبکه‌های عمیق و مشتق‌گیری خودکار، ساختاری عام برای حل مسائل مستقیم و معکوس ارائه دادند که نیازی به تولید شبکه محاسباتی، خطی‌سازی موضعی یا کتابخانه جملات کاندید ندارد. عمومیت این چارچوب و کارایی آن در مسائل متنوع دینامیک سیالات و فیزیک غیرخطی، موج گسترده‌ای از تحقیقات نظری و کاربردی را در قالب مقالات مروری جامع نظیر کومو و همکاران \cite{cuomo2022} به همراه داشته است. پژوهش حاضر بر پایه همین چارچوب شکل گرفته و پیاده‌سازی و ارزیابی تحلیلی آن را در هر دو دسته مسائل مستقیم و معکوس دنبال می‌کند.

\section{جمع‌بندی فصل}
\label{sec:ch2-summary}

در این فصل، مبانی نظری مورد نیاز برای فهم و تحلیل شبکه‌های عصبی آگاه از فیزیک مرور شد: تعاریف و معادلات دیفرانسیل جزئی غیرخطی، روش‌های عددی کلاسیک و چالش‌های محاسباتی آن‌ها، مبانی ساختار و بهینه‌سازی شبکه‌های عصبی عمیق، مبانی مشتق‌گیری خودکار و تاریخچه تحولات یادگیری ماشین فیزیک‌پایه. در فصل بعد، فرمول‌بندی ریاضی و الگوریتمی شبکه‌های آگاه از فیزیک به‌تفصیل تبیین خواهد شد.
